\chapter{الدارات على التسلسل}\label{ch:g4:series-circuits}

في السنة الماضية أظلمت سلسلة أضواء زينة لأن مصباحًا واحدًا فقط
كان ناقصًا --- ووعدنا بأن لترتيبها اسمًا. وها هو ذا: \omterm{def:g4:series-circuits:series}{الدارة على
التسلسل}، أبسط طريقة لوضع لاعبين كثيرين على حلقة واحدة. ولها قواعد
بيت صارمة، \omterm{def:g3:first-electric-circuit:bulb}{والمصباح} اليدوي في درجك يطيعها منذ الأزل.

\section{حلقة واحدة ولاعبون عدّة}

\begin{definition}[على التسلسل]\label{def:g4:series-circuits:series}
يكون اللاعبون في دارة \emph{على التسلسل}\index{دارة على التسلسل}
حين يجلسون على حلقة واحدة، واحدًا بعد الآخر، مثل خرزات في قلادة
واحدة: فطريق التيار يمرّ في الأول، ثم في الثاني، ثم في الثالث،
بلا تفرّع في أي موضع. والدارة المبنيّة هكذا \emph{دارة على
التسلسل}.
\end{definition}

\begin{omfigure}
\begin{tikzpicture}[scale=0.9]
  \draw (0,0) to[battery1, l=بطارية] (0,2.6)
    to[short] (1.6,2.6)
    to[cspst, l=قاطعة] (3.2,2.6)
    to[lamp, l=مصباح 1] (5.2,2.6)
    to[lamp, l=مصباح 2] (7.2,2.6)
    to[short] (7.2,0)
    to[short] (0,0);
\end{tikzpicture}

{\small \omterm{def:g4:series-circuits:series}{دارة على التسلسل}: \omterm{def:g3:first-electric-circuit:battery}{بطارية} وقاطعة ومصباحان على حلقة واحدة
--- وكل لاعب منظوم في القلادة نفسها.}
\end{omfigure}

\begin{method}[كيف تبنيها]\label{met:g4:series-circuits:build}
\omterm{def:g3:first-electric-circuit:battery}{ببطارية} واحدة ومصباحين وقاطعة وأسلاك:
\begin{enumerate}
  \item ابدأ من قطب \omterm{def:g3:first-electric-circuit:battery}{البطارية} $+$ بسلك إلى \omterm{ex:g3:first-electric-circuit:switch}{القاطعة}؛
  \item ثم سلك من \omterm{ex:g3:first-electric-circuit:switch}{القاطعة} إلى \omterm{def:g3:first-electric-circuit:bulb}{المصباح} الأول، ومن \omterm{def:g3:first-electric-circuit:bulb}{المصباح} الأول
        إلى الثاني؛
  \item ثم أغلق الحلقة: من \omterm{def:g3:first-electric-circuit:bulb}{المصباح} الثاني رجوعًا إلى قطب
        \omterm{def:g3:first-electric-circuit:battery}{البطارية} $-$؛
  \item ثم اطقّ \omterm{ex:g3:first-electric-circuit:switch}{القاطعة} وراقب المصباحين معًا.
\end{enumerate}
وتتبّع الحلقة بإصبعك قبل التشغيل: فإن اضطرّ إصبعك يومًا إلى
الاختيار بين طريقين فليست دارةً \omterm{def:g4:series-circuits:series}{على التسلسل}.
\end{method}

\section{قواعد البيت}

\begin{proposition}[قواعد الحلقة الواحدة]\label{prop:g4:series-circuits:rules}
في \omterm{def:g4:series-circuits:series}{الدارة على التسلسل}:
\begin{enumerate}
  \item \emph{الكلّ أو لا شيء}: فراغ واحد في أي موضع --- قاطعة
        مفتوحة، أو \omterm{def:g3:first-electric-circuit:bulb}{مصباح} تالف، أو سلك مرتخٍ --- يوقف كل لاعب في
        الحال؛
  \item \emph{الترتيب لا يهمّ}: بدّل المصباحين، وانقل \omterm{ex:g3:first-electric-circuit:switch}{القاطعة}
        إلى الجهة الأخرى من الحلقة --- يتصرّف كل شيء كما كان
        بالضبط؛
  \item \emph{التقاسم}: كلما كثرت المصابيح على الحلقة خفت توهّج
        كل واحد منها؛ فدفع \omterm{def:g3:first-electric-circuit:battery}{البطارية} يُتقاسم على طول القلادة
        كلها.
\end{enumerate}
\end{proposition}

\begin{example}[رؤية القواعد]\label{ex:g4:series-circuits:seeing}
ابنِ دارة المصباحين واختبر كل قاعدة. فُكّ \emph{أيًّا} من
المصباحين: يعتم الاثنان --- وهي القاعدة الأولى (فالمقبس الفارغ
فراغ). وبدّل المصباحين، أو انقل \omterm{ex:g3:first-electric-circuit:switch}{القاطعة} بينهما: لا تغيّر البتة
--- وهي القاعدة الثانية. والآن أعد البناء \omterm{def:g3:first-electric-circuit:bulb}{بمصباح} واحد: يسطع
وحده سطوعًا أظهر ممّا كان مع رفيقه --- وهي القاعدة الثالثة،
معكوسة.
\end{example}

\begin{example}[أضواء الزينة، وقد حُلّ لغزها]\label{ex:g4:series-circuits:partylights}
اللغز القديم لم يعد لغزًا: فمصابيح الزينة \omterm{def:g4:series-circuits:series}{على التسلسل}، عشرات
الخرزات في قلادة واحدة. \omterm{def:g3:first-electric-circuit:bulb}{والمصباح} الناقص فراغ، والقاعدة الأولى
تُظلم السلسلة كلها. والبحث عن \omterm{def:g3:first-electric-circuit:bulb}{المصباح} المذنب بتجريبها واحدًا
واحدًا علّم أجيالًا من الأسر القاعدةَ الأولى تعليمًا قاسيًا ---
أما السلاسل الحديثة فمبنيّة على نحو آخر، كما ستبيّن السنة
القادمة.
\end{example}

\begin{remark}[القاطعة تأمر الجميع]\label{rem:g4:series-circuits:switch}
للقاعدة الثانية نتيجة مفيدة: أن قاطعة واحدة، موضوعة في أي موضع
من الحلقة، تأمر كل لاعب عليها. ولهذا تستطيع قاطعة جدارية واحدة
أن تتحكّم في صفّ كامل من أضواء السقف --- ولهذا نفسه لا تستطيع
أن تُطفئ مصباحًا \emph{واحدًا} من سلسلة \omterm{def:g4:series-circuits:series}{على التسلسل} من غير أن
تُطفئها كلها.
\end{remark}

\section{بطاريات أكثر في القلادة}

\begin{example}[رأسًا إلى ذيل]\label{ex:g4:series-circuits:batteries}
تستطيع البطاريات أن تصطفّ \omterm{def:g4:series-circuits:series}{على التسلسل} أيضًا --- لكنها أدقّ مزاجًا
من المصابيح: إذ يجب أن تقف \emph{رأسًا إلى ذيل}، بحيث يلمس القطب
$+$ في كل \omterm{def:g3:first-electric-circuit:battery}{بطارية} القطب $-$ في التالية. وبطاريتان مصطفّتان هكذا
تدفعان معًا، فيسطع \omterm{def:g3:first-electric-circuit:bulb}{المصباح} سطوعًا أظهر منه مع واحدة. أدِر إحدى
البطاريتين --- رأسًا إلى رأس --- فيتقاتل الدفعان: فلا يعطي
\omterm{def:g3:first-electric-circuit:bulb}{المصباح} شيئًا.
\end{example}

\begin{omfigure}
\begin{tikzpicture}[scale=0.8]
  \draw (0,0) to[battery1] (0,1.8)
    to[battery1, l=بطاريتان] (0,3.6)
    to[short] (3.6,3.6)
    to[lamp, l=أسطع!] (3.6,0)
    to[short] (0,0);
  \begin{scope}[xshift=6.9cm]
    \draw (0,0) to[battery1, l=بطارية واحدة] (0,1.8)
      to[short] (0,3.6)
      to[short] (3.6,3.6)
      to[lamp, l=أخفت] (3.6,0)
      to[short] (0,0);
  \end{scope}
\end{tikzpicture}

{\small بطاريتان مصطفّتان رأسًا إلى ذيل تدفعان معًا: فيتوهّج
\omterm{def:g3:first-electric-circuit:bulb}{المصباح} نفسه أسطع منه مع \omterm{def:g3:first-electric-circuit:battery}{بطارية} واحدة.}
\end{omfigure}

\begin{example}[داخل المصباح اليدوي]\label{ex:g4:series-circuits:flashlight}
افتح مصباحًا يدويًّا: تنزلق البطاريات فيه واحدة خلف الأخرى في
أنبوب --- وهو اصطفاف رأسًا إلى ذيل جاهز --- ثم يُكمل المصباحُ
وقاطعةُ الزرّ حلقةً واحدة عبر الغلاف. \omterm{def:g3:first-electric-circuit:bulb}{فالمصباح} اليدوي \omterm{def:g4:series-circuits:series}{دارة على
التسلسل} في جيبك. وأنت تعرف الآن لماذا يعتم إذا أُدخلت \omterm{def:g3:first-electric-circuit:battery}{بطارية}
واحدة مقلوبة: رأسًا إلى رأس، فيتلاشى الدفعان.
\end{example}

\section{تمارين}

\begin{exercise}[$\star$]\label{exo:g4:series-circuits:1}
ماذا يعني ``\omterm{def:g4:series-circuits:series}{على التسلسل}''؟ وأيّ صورة من الحياة اليومية يستعملها
الفصل لها؟
\end{exercise}

\begin{exercise}[$\star$]\label{exo:g4:series-circuits:2}
اذكر قواعد بيت \omterm{def:g4:series-circuits:series}{الدارة على التسلسل} الثلاث.
\end{exercise}

\begin{exercise}[$\star$]\label{exo:g4:series-circuits:3}
في دارة المصباحين فُكّ \omterm{def:g3:first-electric-circuit:bulb}{المصباح} 2. فماذا يحدث \omterm{def:g3:first-electric-circuit:bulb}{للمصباح} 1، ولماذا؟
وأيّ قاعدة تعمل هنا؟
\end{exercise}

\begin{exercise}[$\star$]\label{exo:g4:series-circuits:4}
تنقل زوي \omterm{ex:g3:first-electric-circuit:switch}{القاطعة} من قبل \omterm{def:g3:first-electric-circuit:bulb}{المصباح} 1 إلى بعد \omterm{def:g3:first-electric-circuit:bulb}{المصباح} 2. فما الذي
يتغيّر في تصرّف الدارة؟ وأيّ قاعدة تقول ذلك؟
\end{exercise}

\begin{exercise}[$\star$]\label{exo:g4:series-circuits:5}
\omterm{def:g3:first-electric-circuit:battery}{بطارية} واحدة تضيء مصباحًا واحدًا إضاءة جيدة. توقّع السطوع مع:
مصباحين \omterm{def:g4:series-circuits:series}{على التسلسل}؛ \ ثلاثة مصابيح \omterm{def:g4:series-circuits:series}{على التسلسل}. وما الذي
يُتقاسم؟
\end{exercise}

\begin{exercise}[$\star$]\label{exo:g4:series-circuits:6}
كيف يجب أن تقف بطاريتان لتدفعا معًا؟ ارسم الاصطفاف وعلّم $+$
و$-$ على كل \omterm{def:g3:first-electric-circuit:battery}{بطارية}.
\end{exercise}

\begin{exercise}[$\star$]\label{exo:g4:series-circuits:7}
تتبّع حلقة \omterm{def:g3:first-electric-circuit:bulb}{مصباح} يدوي: سمِّ اللاعبين الذين يزورهم التيار من قطب
\omterm{def:g3:first-electric-circuit:battery}{البطارية} $+$ رجوعًا إلى قطبها $-$.
\end{exercise}

\begin{exercise}[$\star$]\label{exo:g4:series-circuits:8}
دارة جرس باب فيها \omterm{def:g3:first-electric-circuit:battery}{بطارية} وزرّ والجرس \omterm{def:g4:series-circuits:series}{على التسلسل}. فأين تستطيع أن
تضيف زرًّا ثانيًا بحيث لا يرنّ الجرس إلا حين يُضغط الزرّان
\emph{كلاهما}؟
\end{exercise}

\begin{exercise}[$\star$]\label{exo:g4:series-circuits:9}
سلسلة من عشرة أضواء زينة \omterm{def:g4:series-circuits:series}{على التسلسل} معتمة. والمصابيح كلها تبدو
سليمة. صِف خطة صبورة، \omterm{def:g3:first-electric-circuit:bulb}{بمصباح} سليم معروف، لإيجاد \omterm{def:g3:first-electric-circuit:bulb}{المصباح}
التالف.
\end{exercise}

\begin{exercise}[$\star\star$]\label{exo:g4:series-circuits:10}
\omterm{def:g3:first-electric-circuit:bulb}{مصباح} توم اليدوي يأخذ بطاريتين ويسطع سطوعًا جميلًا. وبعد تبديل
البطاريتين في العتمة لا يعطي شيئًا --- مع أن البطاريتين
الجديدتين ممتلئتان، \omterm{def:g3:first-electric-circuit:bulb}{والمصباح} سليم، ولا سلك مرتخيًا. فماذا فعل
توم، ولماذا يبقى \omterm{def:g3:first-electric-circuit:bulb}{المصباح} معتمًا؟ (انظر
\cref{ex:g4:series-circuits:batteries}.)
\end{exercise}

\begin{exercise}[$\star\star$]\label{exo:g4:series-circuits:11}
تريد مهندسة أن تضع قاطعة إيقاف طارئ واحدة لآلة يقع محرّكها على
حلقة واحدة مع بطاريته. فتقول: ``أيّ موضع من الحلقة يفي
بالغرض.'' أهي مصيبة؟ ولماذا تفشل خطتها لو كانت في دارة الآلة
طرق متفرّعة؟ (ولفصل السنة القادمة اسم لتلك الطرق.)
\end{exercise}
